\chapter{Induction}
\label{chp.induction}

\minitoc

\lettrine{L}{'un} des outils fondamentaux en logique est l'induction.
Intuitivement, elle peut se voir comme une généralisation du principe de
récurrence. Nous allons cependant adopter un formalisme différent de celui
utilisé pour faire une simple récurrence. Les objets principaux sur lesquels
nous utiliserons l'induction sont les ensembles inductifs et les relations
inductives, que nous présenterons. Ces deux objets sont associés à des
formalismes différents~: le premier aux grammaire en forme de Backus-Naur, et le
deuxième aux points fixes et à la théorie des treillis. Comme la théorie des
treillis sera étudiée plus tard dans cette partie préliminaire, nous ne
traiterons qu'un cas restreint suffisant pour le travail sur l'induction.

L'objectif de ce chapitre est de donner une justification mathématique aux
procédés qui seront utilisés par la suite, et d'offrir un modèle mathématique
derrière le formalisme introduit, pour le lecteur qui en aurait besoin. Le point
essentiel est avant tout de comprendre comment fonctionne une preuve par
induction et comment utiliser les objets inductifs, puisqu'ils seront utilisés
sans arrêt par la suite. Cependant, les justifications mathématiques données
sont partielles~: tous les cas ne se ramènent pas à ceux traités de façon
évidente. Le lecteur le plus prudent devra trouver comment adapter le formalisme
donné dans ce chapitre aux multiples variantes qui seront utilisées sans le dire
par la suite.

On peut trouver plus de détails sur les définitions par induction dans
\cite{winskell1996formal}, au chapitre $3$.

\section{Ensemble inductif}

\subsection{Construction d'un ensemble inductif}

Au niveau intuitif, les ensembles finis semblent avoir une réalité plus
robuste que les ensembles infinis. Il est en effet très facile de se convaincre
à partir de règles simples qu'il existe un ensemble à $3$ éléments, ou à $n$
éléments pour $n$ aussi grand que l'on veut (bien que se convaincre qu'il existe
un ensemble à $300!$ éléments semble légèrement plus long). Cette robustesse
découle du fait qu'on peut explicitement les construire, et cette possibilité
n'existe que pour un ensemble fini. Pourtant, l'ensemble $\bN$ tend aussi
à être plus facilement accepté qu'un ensemble tel que $\bR / \bQ$.
Un point essentiel qui rend le premier ensemble logiquement plus simple que le
deuxième est qu'il est facile à engendrer~: l'ensemble $\bN$ est constitué
de l'élément $0$ et de l'opération $\Ss$ définie par $n \mapsto n + 1$, et tout
autre élément de $\bN$ peut être construit à partir de ces deux éléments.
Sa structure est donc fondamentalement simple, et peut être décrite en des
termes finis.

C'est exactement cette idée de structure générée par des termes finis qui est
formalisée par les ensembles inductifs. Un ensemble inductif va être un ensemble
obtenu par une liste de générateurs, chaque générateur ayant une arité (un
nombre d'objets qu'il prend en entrée). Dans cette définition d'ensemble
inductif, l'exemple canonique est bien sûr $\bN$ lui-même, qu'on peut
définir par~:
\begin{itemize}
\item un constructeur sans argument, $0$
\item un constructeur à un argument, $\Ss$
\end{itemize}

Avant de donner la définition d'ensemble inductif, nous allons donner un
formalisme pour parler des constructeurs.

\begin{definition}[Signature]
  Une signature est un couple $C,\arite$ tel que $\arite : C \to \bN$.
  On appelle $C$ l'ensemble des constructeurs et, pour $c \in C$, $\arite(c)$
  est appelé l'arité de $c$.
\end{definition}

On appelle en général signature la donnée des symboles qui constituent une
structure munie d'une forme de sens (ici leur arité). La présentation qui en est
faite ici est assez aride dans son formalisme~: l'important est d'accepter que
l'on peut spécifier correctement un ensemble inductif par certaines données.
Pour une description plus systématique, mais utilisant le langage catégorique,
on peut par exemple se référer à \cite{JacobsCLTT}.

Une signature sera généralement donnée sous forme dite de Backus-Naur, abrégée en
BNF (\foreignexpr{Backus-Naur Form}). Cette
présentation se décompose de la façon suivante~:
\[a,b,\ldots \inddefeq \mathrm{cas}\;1 \mid \mathrm{cas}\;2 \mid \ldots\]
où $a,b,\ldots$ représentent les éléments que les constructeurs définissent,
et où chaque cas est la définition d'un nouveau constructeur (ou d'une famille
de constructeurs). Par exemple pour le cas de $\bN$, nous avons :
\[n \inddefeq 0 \mid \Ss n\]
Il est fréquent d'employer des variables qui seront quantifiées hors de la
définition à proprement parler. Par exemple, fixons un ensemble $A$, et donnons
la BNF suivante~:
\[\varlist \inddefeq \nil \mid \cons(a,\varlist)\]
où $a \in A$. Cette définition doit se lire comme la signature
$(\{\nil\} \cup A,\arite)$ où $\arite$
est défini par~:
\[\makeFunBis{\arite}{C}{\bN}{\nil}{0}{a (\in A)}{1}\]
Ainsi, plutôt que de considérer un constructeur à deux arguments ayant un
argument extérieur à la signature, on définit formellement un famille de
constructeurs à un seul argument.

Voyons maintenant comment associer à une signature un ensemble généré par les
constructeurs donnés dans la signature. L'ensemble généré doit être un ensemble
$X$ contenant, pour chaque $x_1,\ldots,x_n \in X$ et $c \in C$ d'arité $n$,
l'objet $c(x_1,\ldots,x_n)$, et ne doit contenir que les objets de cette forme.
Nous procédons alors par le bas~: un premier ensemble est construit par
l'ensemble $C_0 = \compr{c \in C}{\arite(c) = 0}$, puis l'ensemble $C_1$ est
construit par $C_1 = C_0 \cup \compr{c(x_1,\ldots,x_n)}{c \in C, x_1,\ldots,
x_n \in C_0, \arite(c) = n}$ et ainsi de suite. Comme $c$ est simplement un
élément dans notre cas, écrire $c(x_1,\ldots,x_n)$ n'a pas de sens, c'est
pourquoi l'on va utiliser à la place $(c,x_1,\ldots,x_n)$.

\begin{definition}[Ensemble inductif sur une signature]
  Soit $(C,\arite)$ une signature, on définit la suite d'ensembles
  $(X_i)_{i\in\bN}$ par :
  \begin{itemize}
  \item $X_0 \defeq \vide$
  \item $X_{n+1} \defeq
    \compr{(c,x_1,\ldots,x_p)}{c\in C, (x_1,\ldots,x_p)\in (X_n)^p,
    \arite(c) = p }$
  \end{itemize}

  L'ensemble inductif engendré par $(C,\alpha)$ est alors l'ensemble
  \[X = \bigcup_{n\in \bN} X_n\]
\end{definition}

La définition donnée n'est pas exactement celle décrite plus haut, mais la
proposition suivante assure que l'union finale génère bien le même ensemble
avec les deux méthodes.

\begin{proposition}
  Soit $(C,\arite)$ une signature, $X$ l'ensemble inductif engendré par cette
  signature et $(X_n)$ la suite précédemment construite. Alors
  \[\forall n,m \in \bN, n \leq m \implies X_n \subseteq X_m\]
\end{proposition}

\begin{proof}
  On procède par récurrence sur $n$~:
  \begin{itemize}
  \item comme $X_0 = \vide$, il est évident que
    $\vide \subseteq X_m$ pour tout $m \in\bN$.
  \item supposons que $X_n \subseteq X_m$ pour tout $m \geq n$. Alors
    \[X_{n+1} = \compr{(c,x_1,\ldots,x_p)}{c\in C, (x_1,\ldots,x_p)\in(X_n)^p,
    \alpha(c) = p}\]
    mais par inclusion, comme $(x_1,\ldots,x_p)\in (X_n)^p$, on en déduit que
    $(x_1,\ldots,x_p)$ est aussi dans $(X_m)^p$, pour tout $m \geq n$. Ainsi
    $(c,x_1,\ldots,x_p) \in X_{m+1}$ pour tout $m \geq n$, donc
    $X_{n+1}\subseteq X_m$ pour tout $m \geq n+1$.
  \end{itemize}
\end{proof}

Le lemme suivant est un outil de base pour étudier des ensembles inductifs.

\begin{lemma}[Lecture unique]\label{lem.lecture_unique}
  Soit une signature $(C,\arite)$ et l'ensemble $X$ engendré par cette
  signature. Alors pour tout élément $x \in X$, il existe $c\in C$ et
  $x_1,\ldots,x_p \in X$ (possiblement une famille vide si $\arite(c) = 0$)
  telle que $p = \arite(c)$ et $x = (c,x_1,\ldots,x_p)$.
\end{lemma}

\begin{proof}
  Soit $(X_n)$ la suite d'ensemble définie précédemment telle que $X$ en est
  l'union. Si $x\in X$, alors on trouve $n \in \bN$ tel que $x\in X_n$.
  Par disjonction de cas sur ce $n$, on prouve le résultat~:
  \begin{itemize}
  \item si $n = 0$ alors l'hypothèse $x \in X_0$ signifie qu'on a
    $x\in \vide$~: par absurdité de la prémisse, la conclusion est vraie.
  \item si $n = m+1$, alors $x\in X_n$ signifie que
    \[x \in \compr{(c,x_1,\ldots,x_p)}{c\in C, (x_1,\ldots,x_p)\in(X_m)^p,
    \alpha(c) = p}\]
    d'où le résultat par définition.\qedhere
  \end{itemize}
\end{proof}

\begin{remark}
  Nous le verrons dans le \cref{chp.alg}, mais une BNF peut aussi se lire comme
  engendrant un langage, c'est-à-dire un ensemble de suites finies sur un
  ensemble fixé à l'avance. On peut par exemple décider de lire l'entier
  $\Ss\Ss\Ss 0$
  comme une simple suite de symboles, ici la suite $(\Ss,\Ss,\Ss,0)$. Dans ce
  cas, le lemme de lecture unique est non trivial et parfois faux, selon la
  forme de
  notre grammaire, et la structure stratifiée d'union dénombrable n'est plus
  inhérente à l'ensemble construit. Cette autre approche, qui est celle utilisée
  dans \cite{cori2003logique}, a le mérite d'être la plus syntaxique possible
  (il n'y a aucune forme d'interprétation des symboles), mais elle nous semble à
  la fois moins systématique et moins intuitive, c'est pourquoi nous lui
  préférons cette structure stratifiée. Cette structure est plus facilement
  décrite de façon catégorique par des $F$-algèbres initiales, comme on peut le
  voir présenté dans \cite{JacobsCLTT} et qui donne un sens plus profond au
  \cref{thm.PU.ind}.
\end{remark}

\subsection{Récursion et induction}

Maintenant qu'une construction a été donnée d'un ensemble inductif, il faut
vérifier que le comportement que l'on a décrit est en accord avec le
comportement réel de l'ensemble que l'on a construit. Nous avons dit que
l'ensemble engendré par $(C,\arite)$ doit contenir exactement les éléments de
la forme $c(x_1,\ldots,x_n)$ où $x_i \in X$ pour tout $i\in\{1,\ldots,n\}$, mais
une autre façon de penser ce fait
%que l'ensemble ne contient que des applications de constructeurs
est qu'une fonction partant d'un ensemble
inductif est exactement spécifiée par son comportement sur les constructeurs.
De telles fonctions sont appelées récursives, car elles peuvent faire appel à
elles-mêmes pour s'appliquer sur les arguments d'un constructeur, comme nous le
verrons en pratique. Nous verrons ensuite que ce principe de définition
récursive peut se modifier pour donner le principe d'induction, un analogue à
la preuve par récurrence pour un ensemble inductif quelconque.

\begin{theorem}[Propriété universelle des ensembles inductifs]\label{thm.PU.ind}
  Soit $(C,\arite)$ une signature, et $X$ l'ensemble associé à cette signature.
  Soit un ensemble $Y$ quelconque.
  Soit une famille de fonctions $\{f_c\}_{c\in C}$ telles que pour tout $c\in C$,
  $f_c : Y^{\arite(c)} \to Y$ (avec la convention que $Y^0 = \{*\}$ est un
  singleton quelconque). Alors il existe une unique fonction $f : X \to Y$
  telle que
  \[\forall c\in C, \forall (x_1,\ldots,x_p)\in X^{\arite(c)},
  f((c,x_1,\ldots,x_p)) = f_c (f(x_1),\ldots,f(x_p))\]
\end{theorem}

On peut représenter l'équation précédente par le diagramme suivant, où l'égalité
signifie que le diagramme commute, c'est-à-dire que les deux chemins possibles
pour aller d'un coin à l'autre du carré sont égaux.

\begin{center}
  \begin{tikzcd}
    X^{\arite(c)} \ar[r,"f^{\arite(c)}"]\ar[d,"c"] & Y^{\arite(c)} \ar[d,"f_c"] \\
    X \ar[r,"f"] & Y
  \end{tikzcd}
\end{center}

\begin{proof}
  Soit $(X_n)$ la suite d'ensembles construite précédemment pour définir $X$. On
  va prouver par récurrence sur $n$ la proposition suivante~:
  \begin{multline*}
    \forall n\in \bN, \exists ! f_n : X_n \to Y, \forall c \in C,\\
    \forall (x_1,\ldots,x_p)\in (X_n)^{\arite(c)}, f_n((c,x_1,\ldots,x_p))
    = f_c(f_n(x_1),\ldots,f_n(x_p))
  \end{multline*}
  \begin{itemize}
  \item Si $n = 0$, il existe une unique fonction $f_0 : \vide \to Y$ et
    elle vérifie la propriété par vacuité.
  \item Soit $n \in \bN$. Supposons qu'il existe une unique fonction
    $f_n : X_n \to Y$ telle que
    \[\forall c\in C, \forall (x_1,\ldots,x_p)\in (X_n)^{\arite(c)},
    f_n((c,x_1,\ldots,x_p)) = f_c(f(x_1),\ldots,f(x_p))\]
    On définit alors
    \[\makeFun{f_{n+1}}{X_{n+1}}{Y}{(c,x_1,\ldots,x_p)}
              {f_c(f_n(x_1),\ldots,f_n(x_p))}\]
    On remarque que cette fonction vérifie bien la propriété. De plus, si une
    autre fonction $g$ vérifie la propriété, alors pour $x\in X_{n+1}$, on
    trouve par
    le \cref{lem.lecture_unique} $c\in C$ et $x_1,\ldots,x_p \in X_n$ tels que
    $x = (c,x_1,\ldots,x_p)$. On a alors
    \begin{align*}
      f_{n+1}(x) &= f_{n+1}((c,x_1,\ldots,x_p)) & \text{par l'égalité sur }x\\
      &= f_c(f_n(x_1),\ldots,f_n(x_n)) & \text{par définition }f_{n+1}\\
      &= g((c,x_1,\ldots,x_p))& \text{par hypothèse sur }g\\
      &= g(x) & \text{par l'égalité sur }x
    \end{align*}
    Donc pour tout $x\in X_{n+1}$, $f_{n+1} = g(x)$, ce qui signifie que
    $f_{n+1} = g$, d'où l'unicité de $f_{n+1}$.
  \end{itemize}

  \noindent
  Soit $x \in X$, par définition on trouve $n\in \mathbb N$ tel que $x \in X_n$,
  et on peut donc définir $f(x) = f_n(x)$. Pour montrer que la fonction est
  unique, il suffit de remarquer que toute fonction $X \to Y$ vérifiant les
  prémisses du théorème induit une fonction sur chaque $X_n$, et doit donc
  coïncider avec chaque $f_n$ sur $X_n$.
\end{proof}

\begin{remark}
  Pour une signature $(C,\arite)$ et un ensemble associé $X$, chaque $c\in C$
  peut maintenant s'interpréter comme une fonction $c : X^{\arite(c)}\to X$. Nous
  confondrons dorénavant le constructeur et la fonction associée, et écrirons
  donc sans distinction $(c,x_1,\ldots,x_p)$ et $c(x_1,\ldots,x_p)$ pour un
  constructeur $c \in C$.
\end{remark}

Cet outil nous permet maintenant de définir des fonctions dont le domaine est
un ensemble inductif en utilisant sa structure. Lorsqu'on définira une fonction
en utilisant le \cref{thm.PU.ind} sur un ensemble inductif $X$, on utilisera
généralement la formulation \og par induction sur la structure de $X$\fg voire
directement \og par induction sur $x$\fg si le contexte rend clair que $x$ est
un élément de $X$.

\begin{example}%\label{ex.double}
  Donnons un premier exemple de fonction récursive~: la fonction
  $\double :n\mapsto 2n$, définie de $\bN$ dans $\bN$. En effet, on peut
  la décrire par
  \begin{itemize}
  \item $\double(0) = 0$
  \item $\double(\Ss n) = \Ss \Ss \double(n)$
  \end{itemize}
  nous donnant alors une définition de $\double$ grâce au théorème précédent.
\end{example}

\begin{example}
  Un exemple à la fois d'ensemble inductif et de fonction récursive est le
  suivant. Soit $A$ un ensemble quelconque, on définit la signature des listes
  sur $A$ par la grammaire suivante~:
  \[\varlist \inddefeq \nil\mid \cons(a,\varlist)\]
  où $a\in A$. L'ensemble $\List(A)$ est alors l'ensemble inductif
  associé. On définit alors la fonction $\ListLen{\tok}$ donnant la longueur
  d'une
  liste~:
  \begin{itemize}
  \item $\ListLen{\nil} = 0$
  \item $\ListLen{\cons(a,\varlist)} = 1 + |\varlist|$
  \end{itemize}
\end{example}

\begin{notation}
  On utilise aussi
  la notation $\Nil$ pour $\nil$ et
  la notation $a \Cons \varlist$ pour $\cons(a,\varlist)$. De plus,
  on considère l'opération $\mathord{\Cons}$ comme associative à droite,
  c'est-à-dire que $a\Cons b \Cons \varlist$ doit se lire comme
  $a\Cons(b\Cons \varlist)$.
\end{notation}

\begin{exercise}
  Soit $A$ un ensemble. Donner une signature définissant l'ensemble
  $\BinTree(A)$ des arbres
  binaires étiquetés par $A$, constitué d'un objet arbre vide $\nil$ et d'un
  constructeur binaire $\node$ prenant en argument un élément $a$ de $A$
  et deux arbres $g$ et $d$ et retournant un nouvel arbre binaire, d'étiquette
  $a$ et dont les deux sous-arbres sont $g$ et $d$.

  Définir une fonction $h : \BinTree(A) \to \bN$ donnant la
  hauteur d'un arbre, c'est-à-dire la longueur du plus long chemin entre la
  racine de l'arbre (la première étiquette) et une feuille (un arbre vide qui est
  sous-arbre). On prend comme convention que $h(\nil) = 0$.

  Définir une fonction $\ListLen{\tok} : \BinTree(A) \to \bN$ donnant le
  nombre d'étiquettes d'un arbre.
\end{exercise}

\begin{exercise}
  En utilisant la structure inductive de $\bN$, définir la fonction
  $\tok+\tok : \bN \times \bN \to \bN$. On pourra pour cela définir
  la fonction $n \mapsto (m \mapsto n + m)$ et faire une fonction récursive sur
  l'argument $n$.
\end{exercise}

Un procédé similaire est celui d'induction. Là où la récursion nous permet de
définir une fonction depuis un ensemble inductif, l'induction va nous permettre
de faire une preuve sur un ensemble inductif. Il s'agit donc, au lieu de donner
une fonction $X \to Y$, de donner un prédicat $P \subseteq X$ et de montrer que
$P = X$.

\begin{theorem}[Principe d'induction]\label{thm.ind.induct}
  Soit $(C,\arite)$ une signature et $X$ l'ensemble inductif associé. Soit
  $P\subseteq X$ un prédicat sur $X$. Si pour tous $c\in C$ et
  $x_1,\ldots,x_p\in X^{\arite(c)}$, la propriété
  \[x_1\in P \text{ et } x_2\in P \text{ et }\ldots \text{ et }x_p\in P \implies
  c(x_1,\ldots,x_p)\in P\] est vérifiée,
  alors $P = X$.
\end{theorem}

\begin{proof}
  Pour prouver ce résultat, il suffit de montrer que pour tout $n\in\bN$,
  $X_n \subseteq P$ pour $(X_n)$ la suite d'ensembles construisant $X$. En
  effectuant une récurrences sur $n$ :
  \begin{itemize}
  \item Si $n = 0$, alors $\vide \subseteq P$.
  \item Supposons que $X_n \subseteq P$ pour $n \in \bN$. Soit
    $x \in X_{n+1}$. Par définition de $X_{n+1}$, on trouve $c \in C$ et
    $x_1,\ldots,x_p \in X_n$ tels que $x = c(x_1,\ldots,x_p)$. Par hypothèse
    de récurrence, on en déduit que pour tout $i\in\{1,\ldots,p\}$, $x_i\in P$.
    Il vient donc, avec l'hypothèse du théorème sur $P$, que
    $c(x_1,\ldots,x_p)\in P$.

    On en déduit que $X_{n+1}\subseteq P$.
  \end{itemize}

  \noindent
  Ainsi, par récurrence, pour tout $n\in \bN$, $X_n \subseteq P$.
  Cela montre alors que
  \[\bigcup_{n \in \bN} X_n \subseteq P\qedhere\]
\end{proof}

On peut affiner ce résultat en distinguant deux sortes de constructeurs~: d'un
côté les constructeurs d'arité $0$, qui sont des constantes, et de l'autre les
constructeurs d'arité supérieure à $1$. Le théorème précédent nous dit que pour
prouver une proposition $P$ sur un ensemble inductif, il suffit de prouver qu'il
contient les constantes et qu'il est stable par chaque constructeur. On remarque
que dans le cas de $\bN$, engendré par $0$ et $\Ss$, ce principe nous dit
qu'une partie $P$ contenant $0$ et telle que
$\forall n\in\bN, n\in P \implies \Ss n\in P$ est exactement $\bN$~:
c'est le principe de récurrence.

Ainsi, puisque nous prouvons le principe d'induction à partir du principe de
récurrence sur $\bN$, et puisque le principe de récurrence sur $\bN$ est un cas
particulier du principe d'induction, les deux sont logiquement équivalents.
L'induction, cependant, est conceptuellement plus intéressante puisqu'elle peut
s'utiliser dans plus de cas.

\begin{exercise}
  On considère $\bN$ comme un ensemble inductif, et la fonction $\double$
  précédemment définie. Montrer que pour tout $n \in \bN$,
  $\double(n)$ est pair. On prendra comme définition de pair le prédicat
  $\pair(n) \defeq \exists m\in\bN, n = 2 \times m$.
\end{exercise}

\begin{exercise}
  Soit un ensemble $A$. On définit
  $\tok \append \tok : \List(A) \times \List(A)\to \List(A)$ par
  induction sur le premier argument~:
  \begin{itemize}
  \item Pour tout $\varlist\in \List(A)$, $\nil \append \varlist = \varlist$.
  \item Pour tous $\varlist,\varlist'\in\List(A), a\in A$,
    $\cons(a,\varlist) \append \varlist' = \cons(a,\varlist \append \varlist')$.
  \end{itemize}

  \noindent
  Montrer que
  \[\forall \varlist,\varlist'\in \List(A),
  |\varlist\append \varlist'|=|\varlist|+|\varlist'|\]
\end{exercise}

\section{Relation inductive}

\subsection{Construction d'une relation inductive}

Une façon de considérer les ensembles inductifs est, étant donnée une signature,
de prendre le plus petit ensemble stable par les constructeurs de cette
signature. Le problème, pour faire cela, est que nous n'avons pas d'ensemble sur
lequel travailler a priori, au sens où l'ensemble inductif que l'on a construit
précédemment n'est pas un élément ou une partie d'un ensemble déjà construit en
amont. C'est pourquoi la construction précédente définissait chaque ensemble
intermédiaire pour en prendre l'union~: l'existence de chaque ensemble de
l'union est à spécifier. Au contraire, pour définir les relations inductives,
nous avons un ensemble ambiant et nous pouvons donc travailler sur la notion de
plus petit ensemble stable puisqu'une relation, au sens d'une partie d'un
ensemble $X$ donné (souvent $X$ est un produit cartésien), vit naturellement
dans $\powerset(X)$. La notion de stabilité mérite un traitement systématique et
une syntaxe appropriée, puisque nous construirons très souvent des relations par
induction. Nous allons donc commencer par donner une description formelle de la
syntaxe des règles d'inférence. Mais avant, rappelons la définition d'une
relation.

\begin{definition}[Relation]
  Soit $X$ un ensemble. On appelle relation sur $X$ une partie $R\subseteq X^n$
  où $n\in\bN$ est appelé l'arité de la relation $R$.
\end{definition}

Une règle d'inférence, elle, va relier des relations.

\begin{definition}[Règle d'inférence]
  On appelle règle d'inférence une présentation de la forme suivante:
  \begin{prooftree}
    \AxiomC{$P_1$}
    \AxiomC{$P_2$}
    \AxiomC{$\cdots$}
    \AxiomC{$P_n$}
    \RightLabel{$\ruleR$}
    \QuaternaryInfC{$P$}
  \end{prooftree}
  Où $P_1,\ldots,P_n$ sont appelées les prémisses de la règle, et $P$ est
  appelée la conclusion de la règle. On dit que la règle est valide si, lorsque
  toutes les prémisses de la règle sont vérifiées, alors la conclusion est aussi
  vérifiée.

  Une règle peut contenir plusieurs paramètres, par exemple
  \begin{prooftree}
    \AxiomC{$a = b$}
    \AxiomC{$P(a)$}
    \BinaryInfC{$P(b)$}
  \end{prooftree}
  auquel cas la règle est valide lorsqu'elle l'est pour chaque instance
  possible de ces paramètres. Dans l'exemple, on suppose que $a$ et $b$ sont
  quantifiés sur un certain ensemble $A$, et la règle est donc valide lorsque
  $\forall a,b\in A, ((a = b) \text{ et } P(a)) \implies P(b)$. On dira que la
  règle est valide aussi lorsque l'on fixe un paramètre, mais on dira alors
  \og la règle $\ruleR$ est valide pour $X = R$\fg{}, dans le cas où l'on
  fixe le paramètre $X$ comme valant $R$. Pour signifier qu'une assignation
  des variables $x_i$ à des objets $a_i$ vérifie un énoncé $P$, on notera
  $(x_i \mapsto a_i) \satisf P$.
\end{definition}

Supposons donnée une règle de la forme
\begin{prooftree}
  \AxiomC{$P_1(R,x_1,\ldots,x_k)$}
  \AxiomC{$P_2(R,x_1,\ldots,x_k)$}
  \AxiomC{$\cdots$}
  \AxiomC{$P_n(R,x_1,\ldots,x_k)$}
  \RightLabel{$\ruleR$}
  \QuaternaryInfC{$P(R,x_1,\ldots,x_k)$}
\end{prooftree}
dont les paramètres sont $R$, une partie de $E^p, p \in \bN$ et les
variables $x_i$, des éléments de $E$. Vérifier que la règle est valide signifie
que pour toute instance des $x_i$ vérifiant tous les $P_i$, la proposition $P$
est valide aussi. En supposant que la conclusion est de la forme
$R(t_1,\ldots,t_p)$, où $t_1,\ldots,t_p$ sont des termes définis en fonction
de $x_1,\ldots,x_k$, alors la validité de $r$ signifie que, si l'on définit
l'ensemble $X$ des $(x_1,\ldots,x_k)$ vérifiant $P_i(R,x_1,\ldots,x_k)$, alors
l'ensemble des termes $t_1,\ldots,t_p$ est inclus dans $R$.

Dans la suite de ce chapitre, on considérera toujours des règles se finissant
par une clause de la forme $R(f_1(x_1,\ldots,x_k),\ldots,f_p(x_1,\ldots,x_k))$.
Une règle $\ruleR$ définit alors la fonction suivante, qui considère
l'ensemble des $R(t_1,\ldots,t_k)$~:
\begin{multline*}
  X \longmapsto\compr{(f_1(a_1,\ldots,a_k),\ldots,f_p(a_1,\ldots,a_k))}{\\
  a_1,\ldots,a_k \in E, \forall i \in \{1,\ldots,n\},
  (R\mapsto X,x_1\mapsto a_1,\ldots,x_k\mapsto a_k) \satisf P_i}
\end{multline*}
On note cette fonction $f_{\ruleR}$. On suppose maintenant une hypothèse qui,
en pratique, sera toujours vérifiée~: la fonction $f_{\ruleR}$ est croissante.
Le cas le plus fréquent est celui où les prémisses de la règle sont directement
des instances de la forme $R(t_1,\ldots,t_n)$, auquel cas il est clair que
notre fonction est croissante.

La validité de $\ruleR$ pour une relation $R$ revient à l'inclusion
$f_{\ruleR}(R)\subseteq R$, que l'on peut aussi exprimer par
$f_{\ruleR}(R)\cup R = R$. Cela nous pousse à définir la fonction suivante, qui
nous sera essentielle pour les prochaines constructions.

\[\makeFun{\ruleR}{\powerset(E^p)}{\powerset(E^p)}{X}{X\cup f_{\ruleR}(X)}\]

Donnons quelques exemples du résultat de $\mathrm r$ pour des règles simples.

\begin{example}
  Prenons la règle
  \begin{prooftree}
    \AxiomC{}
    \RightLabel{$\reflR$}
    \UnaryInfC{$R(x,x)$}
  \end{prooftree}
  exprimant que $R$ est réflexive. La fonction associée à $\reflR$ est
  alors $R \mapsto R \cup \compr{(x,x)}{x \in E}$. Remarquons que l'absence de
  prémisse signifie que la conclusion est toujours valide.
\end{example}

\begin{example}
  La règle précédente n'utilisait pas $R$ dans ses prémisses, un exemple
  l'utilisant est la règle exprimant la transitivité :
  \begin{prooftree}
    \AxiomC{$R(x,y)$}
    \AxiomC{$R(y,z)$}
    \RightLabel{$\transR$}
    \BinaryInfC{$R(x,z)$}
  \end{prooftree}
  Dans ce cas, l'image de $R$ par $\transR$ est l'ensemble
  \[R \cup \compr{ (x,z) \in E^2}{\exists y \in E, R(x,y)\text{ et } R(y,z)}\]
  En effet, même si $y$ n'apparaît pas dans la conclusion, il est nécessaire
  qu'un tel $y$ existe pour que le triplet $(x,y,z)$ vérifie la condition
  de $f_{\transR}$.
\end{example}

Supposons maintenant que nous ayons une relation $R$ et une règle $\ruleR$,
$R$ ne vérifiant pas forcément $\ruleR$. Nous cherchons alors à construire à
partir de $R$ une relation vérifiant $\ruleR$. Une remarque essentielle~: si
$R$ vérifie $\ruleR$, alors $R$ est un point fixe de la fonction $\ruleR$
associée. On va ainsi chercher à créer un point fixe de la fonction associée.
Pour cela, nous allons donner un théorème essentiel en logique~: le théorème de
Knaster-Tarski.

\begin{theorem}[Knaster-Tarski (faible)\cite{Knaster}]\label{thm.ind.KT}
   Soit un ensemble $E$ quelconque, et une fonction
  $f : \powerset(E) \to \powerset (E)$ croissante pour $\subseteq$,
  c'est-à-dire telle que
  \[\forall A,B\in\powerset(E), A\subseteq B \implies f(A)\subseteq f(B)\]
  Alors il existe un plus petit point fixe de $f$ pour $\subseteq$, ou de façon
  équivalente il existe un élément $A\in \powerset(E)$ tel que $f(A)=A$ et tel
  que pour tout $B\in\powerset(E)$ vérifiant $f(B)=B$, on a l'inclusion
  $A\subseteq B$.
\end{theorem}
\begin{proof}
  Pour commencer, définissons l'ensemble des pré-points fixes de $f$, qui est
  \[\prefix(f) = \compr{A \in \powerset(E)}{f(A)\subseteq A}\]
  Cet ensemble est non vide, puisqu'il contient $E$.
  Soit $A = \bigcap \prefix(f)$, montrons que $A$ est dans
  $\prefix(f)$~:

  Pour cela, il suffit de montrer que $f(A)\subseteq A$, mais comme
  $A$ est une intersection, il suffit de montrer que pour tout
  $B\in\prefix(f), f(A)\subseteq B$. Pour montrer cela, on
  remarque par transitivité de $\subseteq$ qu'il suffit de montrer que
  $f(A)\subseteq f(B)$ pour tout $B\in \prefix(f)$, mais cela est
  direct en utilisant le fait que $f$ est croissante et que pour tout
  $B\in \prefix(f)$, $A \subseteq B$. Ainsi
  $f(A)\subseteq A$.

  De plus, comme $f(A)\subseteq A$, on en déduit par croissance de
  $f$ que $f(f(A))\subseteq f(A)$, c'est-à-dire que $f(A)$ est
  lui aussi un élément de $\prefix(f)$~: comme $A$ en est l'intersection,
  cela signifie que $A \subseteq f(A)$, d'où en
  utilisant l'inclusion précédente, $f(A) = A$.

  De plus, si $B$ est un point fixe de $f$, alors $f(B)\subseteq B$ donc par
  définition de $A$, $A \subseteq B$.
\end{proof}

\begin{exercise}
  Soit $E$ un ensemble quelconque, $f : \powerset(E) \to \powerset(E)$ une
  fonction croissante pour l'inclusion et $A \in \powerset(E)$ tel que
  $A\subseteq f(A)$. Montrer que l'on peut étendre le \cref{thm.ind.KT} pour
  trouver un plus petit point fixe $\alpha$ tel que $A\subseteq \alpha$.
\end{exercise}

\begin{exercise}
  Soit $E$ un ensemble quelconque, et $f : \powerset(E) \to \powerset(E)$ une
  fonction croissante pour l'inclusion. On définit
  \[\makeFun{g_f}{\powerset(E)}{\powerset(E)}{X}{X\cup f(X)}\]
  Montrer que $g_f$ est croissante pour l'inclusion et que pour tout
  $X\in\powerset(E), X\subseteq f(X)$.
\end{exercise}

Des deux exercices précédents, on peut déduire le résultat suivant~:

\begin{corollary}[Définition d'une relation inductive]
  Soit une règle $\ruleR$ sur un ensemble $E$ et une relation $R$. Alors il
  existe une plus petite relation $R_{\ruleR}$ contenant $R$ et stable par
  $\ruleR$. On dit
  alors que cette relation $R_{\ruleR}$ est la relation définie par $\ruleR$
  sur $R$. Si
  $R = \vide$, on dira seulement que la relation est définie par $\ruleR$.
\end{corollary}

Cela nous permet alors de définir le principe d'induction sur les relations,
qui est très proche de celui sur les ensembles inductifs.

\begin{theorem}[Induction sur une relation]
  Soit $E$ un ensemble, $R$ une relation $n$-aire sur $E$ et $\ruleR$ une règle
  d'inférence incluant comme paramètre $X$. Soit $P \in \mathcal P(E^n)$
  un prédicat d'arité $n$ sur $E$. Supposons que $R\subseteq P$ et que
  $(X\mapsto P)\satisf \ruleR$. Alors $R_{\ruleR}\subseteq P$.
\end{theorem}

\begin{proof}
  Il suffit de remarquer que $P$ est un point fixe de $\ruleR$ puisque
  $(X\mapsto P)\satisf \ruleR$. Ainsi, comme $P$ est un point fixe de $\ruleR$
  contenant $R$, on en déduit que $P$ contient le plus petit point fixe de
  $\ruleR$ contenant $R$, qui est exactement $R_{\ruleR}$.
\end{proof}

Ce résultat est essentiel pour prouver beaucoup de résultats~: étant donnée une
relation $R$ construite par induction, pour montrer un résultat de la forme
$\forall x_1,\ldots,x_n\in E, R(x_1,\ldots,x_n)\implies P(x_1,\ldots,x_n)$ pour
un certain prédicat $P$, il suffit de montrer que ce prédicat est stable par
la règle définissant $R$.

\begin{exercise}
  Soient désormais $n$ règles $\ruleR_1,\ldots,\ruleR_n$ incluant toutes
  un paramètre $R$ d'arité $p$. Montrer qu'on peut leur associer une fonction
  $\ruleR_{1,\ldots,n} : \mathcal P(E^p) \to \mathcal P(E^p)$ croissante. En
  déduire
  un analogue des propositions précédentes pour un nombre fini de règles.
\end{exercise}

\begin{remark}
  La plupart des relations inductives que nous construirons utilisent plusieurs
  règles, mais l'idée de la construction pour une seule règle suffit. L'exercice
  précédent sert principalement à justifier au lecteur le plus dubitatif que
  le procédé fonctionne effectivement aussi pour plusieurs règles.
\end{remark}

\begin{exercise}
  Soit une relation $R$ et une règle $\ruleR$. On définit la relation $R'$
  comme la relation inductive définie par $\ruleR$ sur $R$, et $R''$ la
  relation inductive définie par la règle $\ruleR$ et
  \begin{prooftree}
    \AxiomC{$R(x_1,\ldots,x_n)$}
    \UnaryInfC{$R''(x_1,\ldots,x_n)$}
  \end{prooftree}
  Montrer que ces deux relations coïncident.
\end{exercise}

\begin{remark}
  On définira dorénavant des relations inductives seulement par des règles.
\end{remark}

\begin{exercise}
  Soit $A$ un ensemble quelconque. On définit sur $\List(A)$ le prédicat
  $\pair$ par induction avec les règles suivantes :
  \begin{center}
    \bottomAlignProof
    \AxiomC{}
    \UnaryInfC{$\pair(\nil)$}
    \DisplayProof
    \qquad
    \bottomAlignProof
    \AxiomC{$\pair(\varlist)$}
    \UnaryInfC{$\pair(a \Cons b \Cons \varlist)$}
    \DisplayProof
  \end{center}
  Montrer l'assertion suivante~:
  \[\forall \varlist \in \List(A), \pair(\varlist) \implies \pair(|\varlist|)\]
  avec le prédicat $\pair$ sur les entiers défini précédemment.
\end{exercise}

\subsection{Dérivation d'arbre}

Une autre utilité de ce formalisme des règles d'inférences est de permettre de
définir une dérivation, qui est une preuve purement syntaxique qu'une certaine
relation est vérifiée.

\begin{definition}[Dérivation]
  Soit une relation $R$ définie par des règles $\ruleR_1,\ldots,\ruleR_n$. Une
  dérivation
  de $R(x_1,\ldots,x_p)$ est un arbre dont la racine est $R(x_1,\ldots,x_p)$ et
  tel que chaque n\oe ud de l'arbre est une instance valide d'une règle parmi
  $\ruleR_1,\ldots,\ruleR_n$.
\end{definition}

\begin{example}
  Soit $A = \{0\}$, dérivons $\pair([0,0,0,0])$ où $[0,0,0,0]$ est une
  abréviation pour
  $0 \Cons 0 \Cons 0 \Cons 0 \Cons \Nil$ :
  \begin{prooftree}
    \AxiomC{}
    \UnaryInfC{$\pair(\Nil)$}
    \UnaryInfC{$\pair([0,0])$}
    \UnaryInfC{$\pair([0,0,0,0])$}
  \end{prooftree}
\end{example}

\begin{remark}
  La notion de dérivation dépend uniquement des règles utilisées~: on n'a donc
  pas besoin de préciser la relation dans laquelle on travaille. Comme on
  travaille en général dans le cas d'une relation définie par les règles, la
  question de si la dérivation dépend des règles ou de la relation en elle-même
  est sans importance.
\end{remark}

Par construction, comme une relation définie par des règles vérifie ces règles,
et étant donnée la définition de \og vérifier une règle\fg{}, il ne fait aucun
doute que si l'on peut dériver $R(x_1,\ldots,x_n)$, alors $R(x_1,\ldots,x_n)$
est vraie.

En logique, nous formalisons les preuves mathématiques comme de tels arbres, car
leur structure simple permet de faciliter l'étude du langage mathématique. Si,
dans la réalité, on n'écrit pas une preuve comme une dérivation (nous le
verrons, cette pratique est beaucoup trop laborieuse), il est communément admis
qu'une preuve satisfaisante peut virtuellement se transformer en un (très grand)
arbre de dérivation.

\'Etudions maintenant la notion de règle admissible et dérivable, correspondant
respectivement à une règle vérifiée par une relation et à une règle dont on peut
syntaxiquement prouver la correspondance avec la règle.

\begin{definition}[Règle admissible]
  Une règle $\ruleR$ est admissible pour une relation $R$ si
  $R\satisf \ruleR$. La règle est admissible pour un ensemble de règles
  $\ruleR_1,\ldots,\ruleR_n$ si la relation définie par ces règles vérifie
  la règle $\ruleR$.
\end{definition}

\begin{definition}[Règle dérivable]
  Soit un ensemble de règles $\ruleR_1,\ldots,\ruleR_n$ et une règle
  $\ruleR$. $\ruleR$ est dite dérivable s'il existe une dérivation de la
  conclusion de $\ruleR$ utilisant (potentiellement) les règles
  $\ruleR_1,\ldots,\ruleR_n$ et (potentiellement) les prémisses de
  $\ruleR$.
\end{definition}

\begin{exercise}
  Montrer qu'une règle dérivable est admissible et que si une règle est
  admissible pour une relation $R$, alors la relation définie par cette règle
  sur $R$ est exactement $R$.
\end{exercise}

\begin{exercise}\label{exo.pair_list}
  Avec les définitions précédentes, montrer que la règle
  \begin{prooftree}
    \AxiomC{$\pair(\varlist)$}
    \UnaryInfC{$\pair([a,b,c,d]\append \varlist)$}
  \end{prooftree}
  est dérivable. Montrer que la règle
  \begin{prooftree}
    \AxiomC{$\pair(a \Cons b \Cons \varlist)$}
    \UnaryInfC{$\pair(\varlist)$}
  \end{prooftree}
  est admissible. Est-elle dérivable~?

  On définit la relation $R$ par les règles suivantes~:
  \begin{center}
    \bottomAlignProof
    \AxiomC{}
    \UnaryInfC{$R(\nil)$}
    \DisplayProof
    \qquad
    \bottomAlignProof
    \AxiomC{$R(\varlist)$}
    \UnaryInfC{$R(a\Cons b \Cons \varlist)$}
    \DisplayProof
    \qquad
    \bottomAlignProof
    \AxiomC{$R(\varlist)$}
    \UnaryInfC{$R(a\Cons b \Cons c \Cons \varlist)$}
    \DisplayProof
  \end{center}
  Montrer que la première règle est valide pour $R$, la seconde l'est-elle~?
\end{exercise}

\begin{exercise}
  Soient des règles $\ruleR_1,\ldots,\ruleR_n$ et une règle $\ruleR$.
  Montrer que toute règle $\ruleR'$ dérivable dans
  $\ruleR_1,\ldots,\ruleR_n$ est encore dérivable dans
  $\ruleR_1,\ldots,\ruleR_n,\ruleR$. En utilisant l'\cref{exo.pair_list},
  est-ce le cas en remplaçant l'hypothèse par $\ruleR'$ admissible ?
\end{exercise}

\begin{remark}
  On a une correspondance entre les relations inductives et les ensembles
  inductifs. Cependant, on remarque qu'on perd de l'information dans le cas des
  relation inductives~: dire simplement $R(x_1,\ldots,x_p)$ ne dit pas comment
  on a dérivé cette instance.
\end{remark}

Enfin, donnons le lemme d'inversion, que l'on peut voir comme un affaiblissement
de l'induction sur une relation, pour une seule étape.

\begin{theorem}[Inversion]
  Soient des règles $\ruleR_1,\ldots,\ruleR_n$, un ensemble $E$ et $R$ la
  relation définie sur $R$ par ces règles. Si $R(x_1,\ldots,x_p)$ pour un tuple
  $(x_1,\ldots,x_p)\in E$, alors il existe $i\in\{1,\ldots,n\}$ et des instances
  $P_1,\ldots,P_k$ des prémisses de $\ruleR_i$ tels que
  \begin{prooftree}
    \AxiomC{$P_1$}
    \AxiomC{$\cdots$}
    \AxiomC{$P_k$}
    \RightLabel{$\ruleR_i$}
    \TrinaryInfC{$R(x_1,\ldots,x_p)$}
  \end{prooftree}
  est vérifiée.
\end{theorem}

\begin{proof}
  Pour prouver ce résultat, il suffit de le prouver par induction sur la
  relation $R$. Mais alors, pour chaque cas de la règle à vérifier, le résultat
  est vrai par hypothèse.
\end{proof}

Ce résultat permet de travailler par disjonction de cas lorsque l'on étudie une
relation inductive. On peut l'assimiler au lemme de lecture unique. Par exemple,
il permet de dire à partir de $\pair(\varlist)$ que $\varlist$ est soit
$\Nil$, soit $a \Cons b \Cons \varlist$ pour deux éléments
$a,b\in A$, en prenant les conventions de l'exercice \ref{exo.pair_list}.

\newpage

\section{Exercices et problèmes}

\subsection{Exercices}

\begin{exercise}[Ordres sur les listes]\label{exo.ord.list}
  Soit $A$ un ensemble. On définit la relation
  $\mathord{\prefOrd}\subseteq\List(A)\times\List(A)$ par les règles
  \begin{center}
    \bottomAlignProof
    \AxiomC{}
    \UnaryInfC{$\varlist \prefOrd \varlist$}
    \DisplayProof
    \qquad
    \bottomAlignProof
    \AxiomC{$\varlist \prefOrd \varlist'$}
    \UnaryInfC{$a \Cons \varlist \prefOrd \varlist'$}
    \DisplayProof
  \end{center}
  Montrer que $\mathord{\prefOrd}$ est une relation d'ordre sur
  $\List(A)$. Est-elle totale~?

  On suppose désormais que $A$ est muni d'une relation d'ordre totale
  $\leq$. On définit la relation lexicographique
  $\mathord{\lexleq}\subseteq\List(A)\times\List(A)$ par les règles
  \begin{center}
    \bottomAlignProof
    \AxiomC{}
    \UnaryInfC{$\nil\lexleq \varlist$}
    \DisplayProof
    \qquad
    \bottomAlignProof
    \AxiomC{$a \leq b$}
    \UnaryInfC{$a \Cons \varlist \lexleq b \Cons \varlist'$}
    \DisplayProof
    \qquad
    \bottomAlignProof
    \AxiomC{$\varlist \lexleq \varlist'$}
    \UnaryInfC{$a \Cons \varlist \lexleq a \Cons \varlist'$}
  \end{center}
  Montrer que $\mathord{\lexleq}$ est une relation d'ordre sur $\List(A)$.
  Est-elle totale~?
\end{exercise}

\begin{exercise}[Clôture transitive]
  Soit $E$ un ensemble et $R$ une relation binaire sur $E$. On définit la
  relation $\tClot R$ par les règles
  \begin{center}
    \AxiomC{$x\mathbin{R}y$}
    \UnaryInfC{$x\mathbin{\tClot R} y$}
    \DisplayProof
    \qquad
    \AxiomC{$x\mathbin{\tClot R}y$}
    \AxiomC{$y \mathbin{R} z$}
    \BinaryInfC{$x\mathbin{\tClot R} z$}
    \DisplayProof
  \end{center}
  Montrer que $\tClot R$ est une relation transitive qui contient $R$. Montrer
  que pour toute relation binaire $R'$ transitive telle que $R\subseteq R'$,
  $\tClot R \subseteq R'$.

  \noindent
  Montrer que pour tous $x,y \in E$, $x\mathbin{\tClot R} y$ si et seulement
  s'il existe une suite finie $x_0,\ldots,x_n$ telle que
  \[\begin{cases}
  x_0 = x \\
  x_n = y \\
  n \neq 0 \\
  \forall i \in \{0,\ldots,n-1\}, x_i R x_{i+1}
  \end{cases}\]
\end{exercise}

\begin{exercise}[Clôture transitive et réflexive]
  Soit $E$ un ensemble et $R$ une relation binaire sur $E$. On définit la
  relation $\rtClot R$ par les règles
  \begin{center}
    \bottomAlignProof
    \AxiomC{}
    \UnaryInfC{$x\mathbin{\rtClot R} y$}
    \DisplayProof
    \qquad
    \bottomAlignProof
    \AxiomC{$x\mathbin{\rtClot R}y$}
    \AxiomC{$y \mathbin{R} z$}
    \BinaryInfC{$x\mathbin{\rtClot R} z$}
    \DisplayProof
  \end{center}
  Montrer que $\rtClot R$ est une relation transitive qui contient $R$. Montrer
  que pour toute relation binaire $R'$ transitive telle que $R\subseteq R'$,
  $\rtClot R \subseteq R'$.

  \noindent
  Montrer que pour tous $x,y \in E$, $x\mathbin{\tClot R} y$ si et seulement
  s'il existe une suite finie $x_0,\ldots,x_n$ telle que
  \[\begin{cases}
  x_0 = x \\
  x_n = y \\
  \forall i \in \{0,\ldots,n-1\}, x_i R x_{i+1}
  \end{cases}\]
\end{exercise}

\begin{exercise}[Clôtures symétriques]
  Soit $E$ un ensemble et $R$ une relation binaire sur $E$. On définit les deux
  relations $\sClotMin R$ et $\sClotMax R$ respectivement par les règles
  \begin{center}
    \bottomAlignProof
    \AxiomC{$x \mathbin{R} y$}
    \UnaryInfC{$x\mathbin{\sClotMax R} y$}
    \DisplayProof
    \qquad
    \bottomAlignProof
    \AxiomC{$y \mathbin{R} x$}
    \UnaryInfC{$x\mathbin{\sClotMax R} y$}
    \DisplayProof
    \qquad
    \text{et}
    \qquad
    \bottomAlignProof
    \AxiomC{$x \mathbin{R} y$}
    \AxiomC{$y \mathbin{R} x$}
    \BinaryInfC{$x \mathbin{\sClotMin R} y$}
    \DisplayProof
  \end{center}
  Montrer que $\sClotMax R$ est une relation symétrique contenant $R$, et que
  pour toute relation $R'$ symétrique contenant $R$, $\sClotMax R \subseteq R'$.

  \noindent
  La relation $\sClotMin R$ est-elle symétrique~? Contient-elle $R$~? Montrer
  que pour toute relation $R'$ symétrique telle que $R' \subseteq R$,
  $R'\subseteq \sClotMin R$.
\end{exercise}

\begin{exercise}[Relation d'équivalence]
  Soit $E$ un ensemble et $R$ une relation binaire sur $E$. On définit
  $\rtsClot R$ par les règles
  \begin{center}
    \bottomAlignProof
    \AxiomC{}
    \UnaryInfC{$x \mathbin{\rtsClot R} y$}
    \DisplayProof
    \qquad
    \bottomAlignProof
    \AxiomC{$x \mathbin{\rtsClot R} y$}
    \AxiomC{$y \mathbin{R} z$}
    \BinaryInfC{$x \mathbin{\rtsClot R} y$}
    \DisplayProof
    \qquad
    \bottomAlignProof
    \AxiomC{$x \mathbin{\rtsClot R} y$}
    \AxiomC{$z \mathbin{R} y$}
    \BinaryInfC{$x \mathbin{\rtsClot R} z$}
    \DisplayProof
  \end{center}
  Montrer que $\rtsClot R$ est une relation d'équivalence, et que pour toute
  relation d'équivalence $\sim$ sur $E$ telle que $R\subseteq\mathord{\sim}$,
  on a l'inclusion $\rtsClot R \subseteq \mathord{\sim}$.

  \noindent
  Montrer que $\rtsClot R = \rtClot{(\sClotMax R)}$. A-t-on toujours l'égalité
  $\rtsClot R = \sClotMax{(\rtClot R)}$~? Si non, donner un contre-exemple.
\end{exercise}

\begin{exercise}[Prédicat sur des listes]
  Soit $A$ un ensemble et $P \in \powerset(A)$. On définit inductivement le
  prédicat $\foraList(P) \in \powerset(\List(A))$ par les règles
  \begin{center}
    \bottomAlignProof
    \AxiomC{}
    \UnaryInfC{$\Nil\in\foraList(P)$}
    \DisplayProof
    \qquad
    \bottomAlignProof
    \AxiomC{$\varlist \in \foraList(P)$}
    \AxiomC{$a \in P$}
    \BinaryInfC{$a\Cons \varlist \in \foraList(P)$}
    \DisplayProof
  \end{center}
  \'Etudions ce prédicat~:
  \begin{enumerate}[label = (\roman*)]
  \item Montrer que $[2,4,6,0]\in\foraList(P)$.
  \item Soit $f : A \to B$ une fonction. On définit la fonction
    $\mapList_f : \List(A) \to \List(B)$ récursivement~:
    \begin{itemize}
    \item $\mapList_f(\Nil) \defeq \Nil$
    \item $\mapList_f(a\Cons \varlist) \defeq f(a) \Cons \mapList_f(\varlist)$
    \end{itemize}
    Soient de plus $P \in \powerset(A), Q \in \powerset(B)$ tels que
    $f(P) \subseteq Q$. Montrer que
    $\mapList_f(\foraList(P)) \subseteq \foraList(Q)$.
  \item On définit récursivement la fonction
    $\filterList_P : \List(A) \to \List(A)$ par~:
    \begin{itemize}
    \item $\filterList_P(\Nil) \defeq \Nil$
    \item $\filterList_P(a \Cons \varlist) \defeq \begin{cases}
      a \Cons \filterList_P(\varlist) \text{ si } a \in P\\
      \filterList_P(\varlist) \text{ si } a \notin P
      \end{cases}$
    \end{itemize}
    Montrer que $\filterList_P(\varlist) \in \foraList(P)$ pour toute liste
    $\varlist \in \List(A)$.
  \end{enumerate}
\end{exercise}

\begin{exercise}[Différents arbres]
  Soit $A$ un ensemble. On définit l'ensemble $\Tree(A)$ des arbres
  $A$-branchants inductivement par
  \[T,T' \inddefeq \nil \mid \node(f)\]
  où $f : A \to \Tree(A)$. Ainsi, un arbre $A$-branchant est soit un arbre vide,
  soit la donnée pour chaque $a \in A$ d'un arbre lui-même $A$-branchant.

  \noindent
  On définit aussi l'ensemble $\PreTree(A)$ des pré-arbres $A$-branchants, par
  \[\PreTree(A) \defeq \compr{T \subseteq \List(A)}{\forall \varlist,\varlist',
  \varlist \prefOrd \varlist' \land \varlist' \in T \implies \varlist \in T}\]
  où $\prefOrd$ est la relation de préfixe donnée dans l'\cref{exo.ord.list}.
  Définir une injection $\Tree(A) \to \PreTree(A)$. Cette injection est-elle
  surjective~?
\end{exercise}

\newpage

\subsection{Problèmes}

\begin{problem}[\'Etude des arbres binaires de recherche]
  A FAIRE
\end{problem}

\begin{problem}[La coinduction]
  L'objectif de ce problème est d'étudier la coinduction. L'induction, présentée
  dans ce chapitre, permet de construire le \og plus petit ensemble défini par
  des règles $\ruleR_1,\ldots,\ruleR_n$\fg. La coinduction, elle, considère le
  \og plus grand ensemble défini par des règles $\ruleR_1,\ldots,\ruleR_n$\fg.

  \pbPart{Construction coinductive}

  \noindent
  On considère une signature $(C,\arite)$ telle que donnée pour un ensemble
  inductif. L'objectif de cette première partie est de construire un ensemble
  similaire à l'ensemble inductif construit par cette signature, mais où les
  objet peuvent être potentiellement infinis. Par exemple, on souhaite que la
  signature
  \[\varlist \inddefeq \Nil \mid a \Cons \varlist\]
  contienne des listes infinies telles que $1\Cons 1 \Cons 1 \Cons\cdots$.
  Pour cela, on va considérer une construction générale appelée la complétion
  idéale.

  \noindent
  Tout d'abord, on considère $X_\bot$, qui va être l'ensemble construit par
  la signature inductive $(C\cup\{\bot\},\arite')$ où
  $\arite'\restr{C} = \arite$ et $\arite'(\bot) = 0$. Ainsi, pour une
  BNF
  \[a \inddefeq e_1\mid \cdots\mid e_n\]
  on associe la nouvelle BNF
  \[a \inddefeq e_1\mid \cdots\mid e_n\mid \bot\]

  \noindent
  On définit la relation $\mathord{<}$ sur $X_\bot$ par
  \begin{center}
    \bottomAlignProof
    \AxiomC{}
    \UnaryInfC{$\bot < c(x_1,\ldots,x_k)$}
    \DisplayProof
    \qquad
    \bottomAlignProof
    \AxiomC{$x < x'$}
    \UnaryInfC{$c(\cdots,x,\cdots) < c(\cdots,x',\cdots)$}
    \DisplayProof
    \qquad
    \bottomAlignProof
    \AxiomC{$x < y$}
    \AxiomC{$y < z$}
    \BinaryInfC{$x < z$}
    \DisplayProof
  \end{center}
  où la seconde règle s'applique à chaque constructeur $c \in C$.

  \begin{pbQuest}
  \item Montrer que $\mathord{<}$ est une relation d'ordre stricte,
    c'est-à-dire qu'elle
    est une relation irréflexive et transitive.
  \item Cette relation est-elle totale~? \textit{Indication~: on pourra
    considérer le cas de la signature des listes.}
  \end{pbQuest}

  \noindent
  Un élément $x \in X_\bot$ est dit total si, en notant $X$ l'ensemble
  construit
  sur la signature $(C,\arite)$, $x \in X$.
  
  \begin{pbQuest}[resume]
  \item Montrer que pour tout $x \in X_\bot$, $x$ est total si et seulement
    s'il
    n'existe pas $y \in X_\bot$ tel que $x < y$.
  \end{pbQuest}

  \noindent
  Une partie $Y \subseteq X_\bot$ est appelée un idéal de $X_\bot$ si
  \begin{itemize}
  \item $Y$ est close par le bas~: si $x,y \in X_\bot$, $x < y$ et $y \in Y$
    alors $x \in Y$
  \item $Y$ est dirigée~: si $x,y \in Y$ alors il existe $z \in Y$ tel
    que $x < z$ et $y < z$
  \item $\bot \in Y$
  \end{itemize}
  
  \noindent
  On note l'ensemble des idéaux de $X_\bot$ par
  \[\ideal(X_\bot) \defeq \compr{I \subseteq X_\bot}{I \text{ est un idéal de }
  X_\bot}\]
  
  \begin{pbQuest}[resume]
  \item Montrer que la fonction
    \[\makeFun{\iota}{X_\bot}{\ideal(X_\bot)}{x}
              {\compr{ y \in X_\bot}{y < x}\cup\{x\}}\]
    est une fonction injective, et que pour tous $x,y \in X_\bot$,
    \[x < y \iff \iota(x) \subsetneq \iota(y)\]
  \item Montrer que si $x \in X_\bot$ est un élément total, alors $\iota(x)$
    est
    maximal dans $\ideal(X_\bot)$~:
    \[\forall I \in \ideal(X_\bot), \iota(x) \subseteq I \implies I =
    \iota(x)\]
  \item Donner un idéal maximal $I \in \ideal(X_\bot)$ qui n'est pas de la
    forme $\iota(x)$ pour $x \in X_\bot$.
  \end{pbQuest}

  \noindent
  On définit désormais
  \[X_\infty \defeq \compr{I \in \ideal(X)}{\forall J \in \ideal(X),
  I \subseteq J \implies I = J}\]
  comme étant l'ensemble coinductif défini par la signature $(C,\arite)$.
  
  \begin{pbQuest}[resume]
  \item Construire une injection $X \to X_\infty$.
  \end{pbQuest}

  \pbPart{L'exemple des flux}

  \noindent
  Soit $A$ un ensemble. On considère la signature $(A,\arite)$ où
  $\arite(a) = 1$ pour tout $a \in A$. On appelle ensemble des flux
  (\foreignexpr{stream} en anglais) sur $A$ l'ensemble coinductif $\Stream(A)$
  associé à cette signature.
  
  \begin{pbQuest}
  \item Donner un élément $1^\infty \in \Stream(\bN)$ représentant la liste
    infinie $1 \Cons 1 \Cons \cdots$.
  \item Quel est l'ensemble inductif associé à cette signature~?
  \end{pbQuest}

  \pbPart{Algèbres et co-algèbres}
  
  \noindent
  Pour exprimer un analogue des \cref{thm.PU.ind} et
  \cref{thm.ind.induct} dans le cas des ensembles coinductifs, on commence par
  étudier une formulation alternative de ces théorèmes offrant une version
  dualte. C'est cette version que nous utiliserons pour la coinduction.
  
  \noindent
  Tout d'abord, reprenons le \cref{thm.PU.ind}. Dans notre formulation,
  construire une fonction $f : X \to Y$ où $X$ est inductif revient à se
  donner,
  pour chaque constructeur $c$, une fonction $f_c : Y^{\arite(c)} \to Y$.
  Une façon alternative de se donner de telles fonctions est de considérer
  d'abord l'ensemble
  \[\sum_{c \in C} X^{\arite(c)} \defeq
  \compr{(c,\vec x)}{\vec x \in X^{\arite(c)}}\]
  
  \begin{pbQuest}
  \item Soit $X$ un ensemble et $(C,\arite)$ une signature. Construire une
    bijection entre l'ensemble des fonctions $\sum_{c \in C} X^{\arite(c)} \to X$
    et des familles $\{f_c\}_{c \in C}$ telles que $f_c : X^{\arite(c)} \to X$.
  \end{pbQuest}
  
  \noindent
  On appelle $(C,\arite)$-algèbre un ensemble $X$ muni d'une fonction
  $f : \sum_{c\in C} X^{\arite(c)} \to X$. Si $(X,f)$ et $(Y,g)$ sont deux
  $(C,\arite)$-algèbres, on appelle morphisme de $(C,\arite)$-algèbres une
  fonction $h : X \to Y$ vérifiant
  \[\forall (c, x_1,\ldots,x_k) \in \sum_{c \in C} X^{\arite(c)},
  h(f(c, x_1,\ldots,x_k)) = g(c,h(x_1),\ldots,h(x_k))\]
  ce que l'on peut représenter par le diagramme commutatif
  \begin{center}
    \begin{tikzcd}
      \sum_{c \in C} X^{\arite(c)} \ar[r,"\sum h"]\ar[d,"f"] &
      \sum_{c \in C} Y^{\arite(c)} \ar[d,"g"] \\
      X \ar[r,"h"] & Y
    \end{tikzcd}
  \end{center}
  
  \noindent
  Une $(C,\arite)$-algèbre $(X,f)$ est dite initiale si pour toute autre
  $(C,\arite)$-algèbre $(Y,g)$ il existe un unique morphisme de
  $(C,\arite)$-algèbre de $(X,f)$ vers $(Y,g)$.
  
  \begin{pbQuest}[resume]
  \item Montrer que si $X$ est l'ensemble inductif associé à une signature
    $(C,\arite)$, alors $X$ muni des constructeurs $c \in C$ est une
    $(C,\arite)$-algèbre.
  \item Montrer que cette $(C,\arite)$-algèbre est initiale.
  \end{pbQuest}

  \noindent
  L'induction permet donc de construire des algèbres initiales, c'est en ce
  sens qu'on peut considérer un ensemble inductif comme \og le plus petit
  possible\fg~: toute autre structure se comportant bien vis à vis de la
  signature donnée contient, en un sens, l'ensemble inductif. La coinduction
  peut se voir comme une notion duale, c'est-à-dire une notion dans lequelle
  on
  inverse le sens de lecture. Ainsi, \og le plus grand possible\fg signifie
  que
  toute autre structure est contenu dans l'ensemble co-inductif, qui est alors
  un objet dit terminal. De même, plutôt que de considérer des algèbres, on
  va considérer des co-algèbres.
  
  \noindent
  On introduit donc la notion de $(C,\arite)$-co-algèbre. Une co-algèbre est
  un ensemble $X$ muni d'une fonction $f : X \to \sum_{c \in C}X^{\arite(c)}$.
  Ainsi, plutôt que la donnée des constructeurs, ce sont les destructeurs qui
  donnent le comportement d'un ensemble co-inductif.
  
  \begin{pbQuest}[resume]
  \item Soit $(C,\arite)$ une signature et $X_\infty$ l'ensemble coinductif qui
    lui est associé. Montrer que chaque constructeur $c \in C$ induit une
    fonction $X_\bot^{\arite(c)} \to X_\bot$ croissante pour $\mathord{<}$.
    En déduire que
    $c$ induit une fonction $X_\infty^{\arite(c)} \to X_\infty$.
  \item Montrer que pour tout $x \in X_\infty$, il existe $c \in C$ tel que
    $x = c(x_1,\ldots,x_{\arite(c)})$. En déduire que $X_\infty$ est une
    $(C,\arite)$-co-algèbre.
  \end{pbQuest}

  \noindent
  Un morphisme de $(C,\arite)$-co-algèbres entre $(X,f)$ et $(Y,g)$ est alors
  la donnée d'une fonction $h : X \to Y$ telle que, en notant la fonction
  \[\makeFun{\sum h}{\sum_{c \in C}X^{\arite(c)}}{\sum_{c \in C}Y^{\arite(c)}}
            {(c,x_1,\ldots,x_{\arite(c)})}{(c,h(x_1),\ldots,h(x_{\arite(c)}))}\]
  on a la propriété
  \[\forall x \in X, \sum h (f(x)) = g(h(x))\]
  ce que l'on peut représenter par le diagramme commutatif
  \begin{center}
    \begin{tikzcd}
      X \ar[r, "h"]\ar[d,"f"] & Y\ar[d,"g"] \\
      \sum_{c \in C} X^{\arite(c)} \ar[r,"\sum h"] &
      \sum_{c \in C} Y^{\arite(c)}
    \end{tikzcd}
  \end{center}
  
  \begin{pbQuest}[resume]
  \item Montrer que $X_\infty$ tel que défini précédemment est une
    $(C,\arite)$-co-algèbre terminale, c'est-à-dire que pour toute autre
    $(C,\arite)$-co-algèbre $(Y,f)$ il existe un unique morphisme de
    co-algèbre de $X_\infty$ dans $(Y,f)$.
  \item Démontrer le théorème de co-récursion pour les ensembles coinductifs~:
    \begin{quotation}
      Soit $(C,\arite)$ une signature et $X_\infty$ l'ensemble coinductif
      associé. Soit $Y$ un ensemble quelconque, $c_Y : Y \to C$ une fonction
      associant à chaque élément de $Y$ un constructeur, et pour tout
      $(y,c_Y(y))$, un tuple $(y_1,\ldots,y_{\arite(c_Y(c))})$.
      Alors il existe une unique fonction $h : Y \to X_\infty$ telle que
      \[\forall y \in Y, h(y) = c_Y(y)(h(y_1),\ldots,h(y_{\arite(c_Y(y))}))\]
%      $c \in C$ et
%      $f : Y \to Y^{\arite(c)}$. Alors il existe une unique fonction
%      $h : Y \to X_\infty$ telle que pour tout $y \in Y$, en notant
%      $y_{f,1},\ldots,y_{f,\arite(c)}$ les coordonnées de $f(y)$, on a
%      \[h(y) = c(h(y_{f,1}),\ldots,h(y_{f,\arite(c)}))\]
    \end{quotation}
  \end{pbQuest}

  Montrons maintenant le principe de co-induction, c'est-à-dire la version
  duale du \cref{thm.ind.induct}. Dans l'induction, le principe est de
  considérer un constructeur $c$ d'arité $n$ et de montrer
  \begin{quotation}
    Pour tous $x_1,\ldots,x_n$, si $P(x_1),\ldots,P(x_n)$ alors
    $P(c(x_1,\ldots,x_n))$.
  \end{quotation}
  Dans le cas co-inductif, on cherche simplement à montrer la propriété
  \og dans l'autre sens\fg~:
  \begin{quotation}
    Pour tous $x_1,\ldots,x_n$, si $P(c(x_1,\ldots,x_n))$ alors il existe
    $i$ tel que $P(x_i)$.
  \end{quotation}
  Pour formaliser correctement cette idée, on commence par définir la notion
  de prédicat productif~: pour un ensemble coinductif $X_\infty$ tel
  qu'introduit précédemment, un prédicat $P\subseteq X_\infty$ est dit
  productif s'il vérifie l'énoncé précédent, c'est-à-dire que pour tout
  constructeur $c$ et éléments $x_1,\ldots,x_n$, si
  $P(c(x_1,\ldots,x_n))$ alors il existe $i$ tel que $P(x_i)$.
  \begin{pbQuest}[resume]
  \item Soit $P$ productif sur $X_\infty$, on définit $\tilde P$ sur $X_\bot$
    par
    \[ \tilde P(x) \defeq \exists x_0 \in X_\infty, P(x_0) \land x \in x_0\]
    Montrer que $\tilde P$ est un prédicat inductif, c'est-à-dire que si
    $\tilde P(x_1),\ldots,\tilde P(x_n)$ et $c$ est un constructeur d'arité
    $n$, alors $\tilde P(c(x_1,\ldots,x_n))$.
  \item En déduire que si $P$ est productif, alors
    \[\forall x \in X_\infty, P(x)\]
  \end{pbQuest}
  
  \pbPart{Application aux flux}

  \begin{pbQuest}
  \item En utilisant les résultats précédents, construire par co-récursion une
    fonction
    \[\makeFunName{\fromF}{\bN}{\Stream(\bN)}\]
    qui associe à $n$ le flux $n \Cons (n+1)\Cons (n+2) \Cons \cdots$.
  \item Soit $A$ un ensemble. Définir par co-récursion une fonction
    \[\makeFunName{\iota}{A^\bN}{\Stream(A)}\]
    qui associe à $(u_n)_{n \in \bN}$ le flux $u_0 \Cons u_1 \Cons \cdots$
  \item Montrer que $\iota$ est une bijection.
  \end{pbQuest}

  \pbPart{Relation coinductive}

  \noindent
  Nous avons étudié les ensembles définis coinductivement. Ces constructions
  sont techniques, car il est difficile de construire correctement un objet
  existant \latinexpr{a priori}. C'est aussi pour cette raison que nous avons
  introduit dans notre construction l'ensemble $\powerset(X_\bot)$, permettant
  d'avoir un surensemble dans lequel travailler.

  Dans le cas des relations définies par induction, le travail est plus simple~:
  on travail \latinexpr{a priori} dans un ensemble de parties, et le théorème
  de Knaster-Tarski invite directement à produire un énoncé dual.

  \begin{pbQuest}
  \item Montrer la version duale du théorème de point fixe de Knaster-Tarski~:
    \begin{quotation}
      Soit $E$ un ensemble quelconque, et $f : \powerset(E) \to \powerset(E)$
      une fonction croissante pour $\subseteq$. Il existe un plus grand point
      fixe de $f$ pour $\subseteq$, c'est-à-dire une partie $A$ telle que
      $f(A) = A$ et pour toute partie $B$ vérifiant $f(B) = B$, on a l'inclusion
      $B \subseteq A$.
    \end{quotation}
  \item En déduire une construction, pour des règles d'inférences telles
    qu'introduites plus tôt dans le chapitre, d'une plus grande relation
    satisfaisant ces règles.
  \end{pbQuest}

  \noindent
  L'induction sur une relation inductive peut aussi être effectuée de façon
  duale. Pour cela, considérons une règle d'inférence
  \begin{center}
    \AxiomC{$P_1$}
    \AxiomC{$\cdots$}
    \AxiomC{$P_n$}
    \RightLabel{$\ruleR$}
    \TrinaryInfC{$P$}
    \DisplayProof
  \end{center}
  L'induction nous dit qu'une relation satisfaisant cette règle contient la
  relation inductivement engendrée par cette règle. Par le même principe, une
  telle relation est contenue dans la relation coinductivement engendrée
  par cette règle. Ainsi, une coinduction sur une relation coinductive demande
  de construire pour chaque élément une relation stable par une règle.

  \noindent
  Pour ce faire, introduisons la notion de co-satisfaction d'une règle. Avec les
  notations précédente, la règle $\ruleR$ est dite co-satisfaite pour
  l'instanciation $X \mapsto R, x_1 \mapsto a_1,\ldots, x_n \mapsto a_n$
  lorsque, si la conclusion $P$ de $\ruleR$ sur cette instanciation des
  paramètres,
  alors il existe une des prémisses $P_i$ de $\ruleR$ qui est vérifiée. On
  notera
  $(X\mapsto R, x_1 \mapsto a_1,\ldots, x_n \mapsto a_n)\cosatisf \ruleR$.

  \begin{pbQuest}[resume]
  \item Soit $E$ un ensemble et
    $\ruleR$ une règle d'inférence telle que présentée ci-dessus,
    à paramètres $X \in \powerset(E^n), x_1,\ldots,x_n \in E$. Soit
    $R$ une relation $n$-aire sur $E$.
    Si $(X \mapsto R, x_1\mapsto a_1,\ldots x_n \mapsto a_n) \cosatisf \ruleR$
    alors la relation inductive engendrée par $\ruleR$ sur la relation
    $R$ est contenue dans la relation coinductivement engendrée par $\ruleR$.
  \item En déduire le théorème de coinduction suivant~:
    \begin{quotation}
      Soit $E$ un ensemble et $\{\ruleR_i\}_{i \in I}$ un ensemble de règles
      d'inférence à paramètre $X \in \powerset(E^n), x_1,\ldots,x_n \in E$.
      Soit $R \in \powerset(E^n)$ une relation $n$-aire sur $E$.
      Si pour tous $(a_1,\ldots,a_n)$ il existe $i \in I$ tel que
      $(X \mapsto R, x_1\mapsto a_1,\ldots,x_n\mapsto a_n)\cosatisf \ruleR_i$
      alors $R$ est incluse dans la relation coinductivement engendrée par
      $\{\ruleR_i\}_{i \in I}$.
    \end{quotation}
  \end{pbQuest}
\end{problem}
